\documentclass[11pt,a4paper]{article}
\usepackage[margin=25mm]{geometry}
\usepackage[T1]{fontenc}
\usepackage{lmodern,microtype}
\usepackage{amsmath,amssymb,amsthm,mathtools,booktabs}
\usepackage{xcolor}
\usepackage{hyperref}
\hypersetup{colorlinks=true,linkcolor=black,citecolor=blue!55!black,
 urlcolor=blue!55!black,pdfauthor={Tatsuhiko Takimoto},
 pdftitle={Sharp first-order energy correction under an Eulerian tidal bound}}
\newtheorem{theorem}{Theorem}
\newtheorem{lemma}[theorem]{Lemma}
\newtheorem{proposition}[theorem]{Proposition}
\newcommand{\R}{\mathbb R}
\newcommand{\dd}{\,\mathrm d}
\newcommand{\opnorm}[1]{\left\lVert #1\right\rVert_2}
\newcommand{\tr}{\operatorname{tr}}
\newcommand{\TV}{\operatorname{TV}}
\numberwithin{equation}{section}
\title{Sharp first-order energy correction under an Eulerian tidal bound for finite-wall Alcubierre profiles}
\author{Tatsuhiko Takimoto}
\date{27 October 2026 \quad (version 1.2)}
\begin{document}
\maketitle

\begin{abstract}
We consider the stationary Alcubierre metric with speed \(v=1/10\) and globally
\(C^2\), monotone radial profiles that match constant interior and exterior
values at radii \(1\) and \(R\in[3/2,3]\). For the Eulerian normal, we bound the
spectral norm of the electric part of the full Riemann tensor by \(\eta v^2\)
at every point of the wall. If \(J_R^*(\eta)\) denotes the infimum of
\(\int_1^R r^2 f'(r)^2\dd r\) over this class, then
\[
 \eta\left[J_R^*(\eta)-\frac{R}{R-1}\right]
 \longrightarrow \frac{1}{3R(R-1)^3}
\]
uniformly in \(R\) as \(\eta\to\infty\). An axial necessary condition gives a
lower bound for every admissible profile. A normalized, smoothed ramp
construction gives an upper bound satisfying the tidal constraint at all
angles. The errors in the scaled lower and upper bounds are respectively
\(O(\eta^{-1/2})\) and \(O(\eta^{-1})\), uniformly in \(R\).
No minimizing profile or uniform higher-derivative bound is assumed.
Eighteen exact rational enclosures illustrate the bounds.
\end{abstract}

\section{Introduction}

Alcubierre's metric provides a concrete setting in which a prescribed
spacetime geometry determines an Eulerian energy density
\cite{alcubierre}. Pfenning and Ford related curvature scales, wall thickness
and quantum-inequality restrictions \cite{pfenningford}. Lobo and Visser
used a volume integral to quantify negative energy and distinguished that
quantity from ADM mass \cite{lobovisser}. The present problem concerns the
cost of imposing a pointwise tidal bound on a fixed finite wall.

Shape optimization is already part of this literature. Bobrick and Martire
\cite[arXiv v2, App.~A.3]{bobrickmartire} vary the radial energy functional
and obtain a harmonic profile. Nat\'ario's divergence-free
shift construction \cite{natario} leads to a different admissible class.
For that class, Bol\'ivar, Vasilev and Abell\'an report a sharp Eulerian
energy minimum and a thin-wall asymptotic through a Stokes-flow
formulation \cite{bolivar}. Here the wall width stays in \([1/2,2]\).
The small scale instead arises from endpoint layers as the allowed
electric tidal norm increases.

General shift constructions by Lentz \cite{lentz} and Fell and Heisenberg
\cite{fellheisenberg}, and the curvature and observer analysis of Santiago,
Schuster and Visser \cite{santiago}, provide further context. In particular,
Eulerian density and energy conditions involving other observers are
different questions. Warp Factory \cite{warpfactory} provides computational
tools for geometry and energy-condition analysis; Le reports
observer-robust verification methods in Warpax \cite{warpax}. Our conclusion is an analytic asymptotic for a
specified scalar profile class. Higher-order variational expansions have
a general framework \cite{braides}; the proof here uses direct inequalities
and an explicit recovery family.

The quantity determined below is the first correction to the harmonic
energy floor under an all-angle electric-Riemann norm constraint.
The outer endpoint forces a linear loss in the rescaled profile.
At the inner endpoint, the vanishing shift permits a much thinner
transition. This asymmetry determines the coefficient.

\section{The variational problem and the main result}\label{sec:problem}

Work in units \(G=c_{\mathrm{light}}=1\), fix \(v=1/10\), and write
\begin{equation}\label{eq:metric}
 \dd s^2=-\dd t^2+\dd x^2+\dd y^2+
 \bigl(\dd z+v(1-f(r))\dd t\bigr)^2,\qquad
 r=(x^2+y^2+z^2)^{1/2}.
\end{equation}
For \(R\in[3/2,3]\), the profile \(f\) is globally \(C^2\) and satisfies
\[
 f=1\ \text{on }[0,1],\qquad f=0\ \text{on }[R,\infty),
 \qquad f'\le0.
\]
Thus \(f'=f''=0\) at both matching radii. There is no uniform bound on
derivatives of order three or higher. Set
\[
 h=1-f,\qquad p=h'=-f',\qquad
 n=\partial_t-vh\,\partial_z.
\]
The spatial vectors \(\partial_x,\partial_y,\partial_z\) form an
orthonormal frame perpendicular to \(n\). Let
\begin{equation}\label{eq:Edef}
 E_{ij}=R_{i\alpha j\beta}n^\alpha n^\beta,\qquad
 \opnorm{E/v^2}\le\eta \quad\text{at every wall radius and angle}.
\end{equation}
Here \(E\) is the electric part of the full Riemann tensor, and
\(\opnorm{\cdot}\) is the Euclidean operator norm. An overall reversal of
the curvature sign does not affect the constraint. The unit lapse makes
the Eulerian normal geodesic, so \(E\) is its geodesic-deviation tensor.
For any fixed nonzero \(v\), both \(E/v^2\) and \(J\) below are independent
of \(v\); the choice \(v=1/10\) fixes the energy conversion
in \eqref{eq:physical-energy}.

Let \(\mathcal A_R(\eta)\) be the class just defined, for \(\eta\ge128\), and
put
\begin{equation}\label{eq:energy-index}
 J(f)=\int_1^R r^2p(r)^2\dd r,\qquad
 J_R^*(\eta)=\inf_{f\in\mathcal A_R(\eta)}J(f).
\end{equation}
For later use, define
\begin{equation}\label{eq:constants}
 C=\frac{R}{R-1},\qquad a=\frac{1}{R(R-1)},\qquad
 K_R=\frac{R^2a^3}{3}=\frac{1}{3R(R-1)^3}.
\end{equation}
In this range \(3/2\le C\le3\), \(1/6\le a\le4/3\), and \(K_R\le16/9\).

\begin{theorem}\label{thm:main}
For every \(R\in[3/2,3]\) and \(\eta\ge128\), the class
\(\mathcal A_R(\eta)\) is nonempty and
\begin{equation}\label{eq:finite-lower}
 \eta\bigl(J_R^*(\eta)-C\bigr)\ge L_R(\eta):=
 \frac{(R-a/\eta)^2a^3}{3}
       \left(1-2\sqrt{a/\eta}\right).
\end{equation}
The constants \(B_-=5\), \(B_+=61000\), and \(\eta_0=128\), independent
of \(R\), satisfy, for all \(\eta\ge\eta_0\),
\begin{equation}\label{eq:error-bounds}
 K_R-B_-\eta^{-1/2}
 \le \eta\bigl(J_R^*(\eta)-C\bigr)
 \le K_R+B_+\eta^{-1}.
\end{equation}
Consequently,
\begin{equation}\label{eq:uniform-limit}
 \lim_{\eta\to\infty}
 \sup_{R\in[3/2,3]}
 \left|\eta\left[J_R^*(\eta)-\frac{R}{R-1}\right]
             -\frac{1}{3R(R-1)^3}\right|=0 .
\end{equation}
No attainment of the infimum is required.
\end{theorem}

The coefficient is sharp because a lower bound for the entire admissible
class and an admissible family have the same first-order limit.
The rates in \eqref{eq:error-bounds} are estimates for the bounds; we do
not assert that they are optimal convergence rates.

\section{Curvature and the energy functional}\label{sec:curvature}

Use the convention
\[
 R^a{}_{cmn}=\partial_m\Gamma^a_{nc}-\partial_n\Gamma^a_{mc}
       +\Gamma^a_{mk}\Gamma^k_{nc}-\Gamma^a_{nk}\Gamma^k_{mc}.
\]
For a stationary shift \(b\,\partial_z\), set \(d=\nabla b\),
\(H=\nabla^2b\) and \(e=(0,0,1)^T\). The nontrivial inverse metric
components are \(g^{tt}=-1\), \(g^{tz}=b\), and \(g^{zz}=1-b^2\).
Direct evaluation of the Christoffel symbols and contraction with
\(n=(1,0,0,-b)\) give
\begin{equation}\label{eq:Egeneral}
 \begin{split}
 E={}&-\frac b2\bigl((He)e^T+e(He)^T\bigr)
       -\frac34dd^T\\
    &-\frac{b_z}{4}(de^T+ed^T)
       +\frac{|d|^2}{4}ee^T.
 \end{split}
\end{equation}
This is a specialization of the flat-slice curvature identities discussed
in \cite[arXiv v2, Sec.~3.1]{santiago}. For example, it also follows by taking
\(\mathsf K=(de^T+ed^T)/2\) and expanding
\[
 E=-\mathcal L_{be}\mathsf K+\mathsf K^2,\qquad
 \mathcal L_{be}\mathsf K=b\,\partial_z\mathsf K+
                  d(e^T\mathsf K)+(\mathsf K e)d^T
\]
in the present stationary convention.

Substitute \(b=vh\). Rotational symmetry about the \(z\)-axis allows the
radial unit vector to be written as \((s,0,u)\), with
\(u\in[-1,1]\) and \(s=\sqrt{1-u^2}\). Then
\begin{equation}\label{eq:radial-E}
 E/v^2=\begin{pmatrix}A&0&B\\0&0&0\\B&0&D\end{pmatrix},
 \quad
 \begin{aligned}
 A&=-\tfrac34p^2(1-u^2),\\
 B&=-\tfrac h2(p'-p/r)su-p^2su,\\
 D&=-hp'u^2-\tfrac{hp}{r}(1-u^2)
                         +\tfrac{1-5u^2}{4}p^2 .
 \end{aligned}
\end{equation}
On the axis, \(E_{zz}/v^2=-hp'-p^2=-(hp)'\). Thus every admissible
profile satisfies
\begin{equation}\label{eq:necessary}
 |(hp)'|\le\eta .
\end{equation}
For the upper construction we need an estimate valid at every angle:
\begin{equation}\label{eq:sufficient}
 \opnorm{E/v^2}\le |hp'|+\frac{hp}{r}+\frac32p^2.
\end{equation}
Here \(h,p\ge0\). The coefficient of \(hp'\) has norm
\((u^2+|u|)/2\le1\). The coefficient of \(hp/r\) has norm
\((\sqrt{1-u^2}+1-u^2)/2\le1\). The remaining terms in
\eqref{eq:radial-E} have total norm at most
\(p^2[3/4+|u|(1+|u|)/4+1/4]\le3p^2/2\).
The axial condition is used only as a necessary condition; it is not
substituted for the all-angle constraint.

The Einstein contraction, with \(G_{ab}=8\pi T_{ab}\), is
\[
 G(n,n)=-\frac{b_x^2+b_y^2}{4},\qquad
 \rho=T(n,n)=-\frac{v^2p^2\sin^2\theta}{32\pi}.
\]
The Hamiltonian expression
\(\bigl((\tr\mathsf K)^2-\tr(\mathsf K^2)\bigr)/2\)
gives the same contraction. Since the spatial slice is Euclidean,
\begin{equation}\label{eq:physical-energy}
 \mathcal E:=\int_{\Sigma_t}\rho\dd V
       =-\frac{v^2}{12}\int_1^R r^2p^2\dd r
       =-\frac{J}{1200}.
\end{equation}
Thus \(J\) is proportional to the magnitude of this negative Eulerian
volume integral. Its interpretation as a slice integral follows
\cite{alcubierre,lobovisser}; it is not an ADM mass.

\section{A lower bound for every admissible profile}\label{sec:lower}

Let \(p_0(r)=C/r^2\). We have \(\int_1^R p_0\dd r=1\) and
\(\int_1^R r^2p_0^2\dd r=C\). Since every profile also has
\(\int_1^R p\dd r=1\), expansion of the square gives the exact identity
\begin{equation}\label{eq:square}
 J-C=\int_1^R r^2(p-p_0)^2\dd r\ge0.
\end{equation}
The reference function \(p_0\) does not satisfy the endpoint conditions.

Put \(w=hp\). Equation \eqref{eq:necessary} and \(w(R)=0\) imply
\begin{equation}\label{eq:outer-cap}
 hp(r)\le\eta(R-r).
\end{equation}
First consider profiles for which \(J\le C+4/\eta<4\).
Since \(a/\eta\le1/96<1/2\le R-1\), the interval
\([R-a/\eta,R]\) lies inside the wall. On it,
\(\int_r^R t^{-2}\dd t=(R-r)/(rR)\le R-r\), so Cauchy--Schwarz gives
\[
 1-h(r)=\int_r^R p(t)\dd t
 \le\left(J\int_r^R t^{-2}\dd t\right)^{1/2}
 \le2\sqrt{R-r}\le2\sqrt{a/\eta}=:\beta<1 .
\]
In fact \(\beta\le2\sqrt{(4/3)/128}<0.21\). On the smaller interval
\(0\le t=R-r\le a(1-\beta)/\eta\), the ordering
\[
 p(r)\le\frac{\eta t}{1-\beta}\le a\le p_0(r)
\]
permits squaring \(p_0-p\ge a-\eta t/(1-\beta)\ge0\).
Restricting \eqref{eq:square} to this interval therefore yields
\[
 \begin{split}
 \eta(J-C)
 &\ge \eta(R-a/\eta)^2
       \int_0^{a(1-\beta)/\eta}
          \left(a-\frac{\eta t}{1-\beta}\right)^2\dd t\\
 &=\frac{(R-a/\eta)^2a^3(1-\beta)}{3}=L_R(\eta).
 \end{split}
\]
The integration interval lies inside the wall. If \(J>C+4/\eta\),
then \(\eta(J-C)>4>K_R\ge L_R(\eta)\), so the same conclusion holds.
Taking the infimum proves \eqref{eq:finite-lower}. The bounds on
\(R\) and \(a\) imply
\begin{equation}\label{eq:lower-asymptotic}
 L_R(\eta)=K_R+O(\eta^{-1/2})
\end{equation}
uniformly in \(R\). More explicitly,
\[
 0\le K_R-L_R(\eta)
 \le K_R\left(\frac{2a}{R\eta}+2\sqrt{a/\eta}\right)
 \le5\eta^{-1/2}.
\]
For the last step use \(K_R\le16/9\), \(2a/R\le16/9\),
\(1/\sqrt{128}<1/11\), and \(4/\sqrt3<7/3\).
This argument applies to all admissible \(C^2\)
profiles, including those far from the energy floor.

\section{An admissible recovery family}\label{sec:recovery}

\subsection{Construction and normalization}

For \(\eta\ge128\), set
\[
 \gamma=\eta-64\ge64,\qquad \delta=\eta^{-3},\qquad
 \epsilon=\eta^{-4}<\delta .
\]
Define on \([1,R]\)
\begin{equation}\label{eq:raw-profile}
 q_\eta(r)=\min\left\{\frac C{r^2},
             \frac C\delta(r-1-\delta)_+,
             \gamma(R-\delta-r)_+\right\}
\end{equation}
and extend it by zero outside the wall. The three active pieces are an
increasing line, \(C/r^2\), and a decreasing line. Their junctions satisfy
\begin{equation}\label{eq:junctions}
 \begin{gathered}
 (\alpha-1-\delta)\alpha^2=\delta,\qquad
          1+\delta\le\alpha\le1+2\delta,\\
 (R-\delta-\zeta)\zeta^2=C/\gamma,\qquad
          R-\delta-C/\gamma\le\zeta\le R-\delta.
 \end{gathered}
\end{equation}
The two intervals are disjoint. The inner line increases while
\(C/r^2\) decreases. The function \((R-\delta-r)r^2\) is unimodal,
is greater than \(C/\gamma\) at \(r=1\), and is decreasing on the
outer junction interval, because
\[
 \frac C\gamma\le\frac3{64}<\frac{R-\delta}{3},\qquad
 R-1-\delta\ge\frac12-128^{-3}>\frac3{64}.
\]
Also \((r-1-\delta)r^2\) is strictly increasing on the inner interval.
These inequalities give the unique stated junctions and decomposition.
In particular,
\[
 0\le q_\eta\le C,\qquad \TV(q_\eta)\le2C,\qquad
       -\gamma\le q_\eta'\le C/\delta\quad\text{a.e.}
\]
Here the derivative on the harmonic piece is at least \(-6\ge-\gamma\).

Write \(I=\int_1^R q_\eta\dd r\).
The missing mass at the inner end is at most \(2C\delta\).
At the outer end, the zero strip loses at most \(C\delta\);
comparison with a ramp under the constant \(C\) bounds the remaining
loss by \(C^2/(2\gamma)\). Hence
\begin{equation}\label{eq:Ibound}
 1-I\le3C\delta+\frac{C^2}{2\gamma}\le\frac8\gamma,
 \qquad \frac78\le I\le1 .
\end{equation}
The last estimate follows from \(C\le3\) and
\(9\eta^{-3}\le 7/[2(\eta-64)]\).

Use the nonnegative even kernel
\[
 \kappa(t)=\frac{15}{16}(1-t^2)^2\mathbf1_{\{|t|\le1\}},
 \qquad \kappa_\epsilon(t)=\epsilon^{-1}\kappa(t/\epsilon).
\]
It has mass one and first absolute moment \(5/16\). Define
\begin{equation}\label{eq:smoothed-profile}
 p_\eta=\frac{q_\eta*\kappa_\epsilon}{I},\qquad
 h_\eta(r)=\int_1^r p_\eta(t)\dd t,\qquad f_\eta=1-h_\eta.
\end{equation}
Extend \(h_\eta\) by 0 and 1 on the inner and outer constant regions.
Since \(\epsilon<\delta\), the support of \(p_\eta\) stays strictly
inside the wall. Since \(q_\eta\) is Lipschitz,
\((q_\eta*\kappa_\epsilon)'=q_\eta'*\kappa_\epsilon\) is continuous;
nonnegativity and unit mass of the kernel preserve the a.e. derivative
bounds. Hence \(f_\eta\) is globally \(C^2\), is monotone, and has both required zero derivatives
at the endpoints. Its normalization is exact.

\subsection{The tidal bound at all wall points}

Put \(A=C/I\le24/7<4\). Then
\[
 p_\eta\le A,\qquad -\gamma/I\le p_\eta'\le A/\delta.
\]
The derivative \(q_\eta'\) is positive only on \((1+\delta,\alpha)\)
and is nonpositive a.e. on \([\alpha,R]\).
If \(p_\eta'(r)>0\), the convolution window meets
\((1+\delta,\alpha)\), so \(r\le1+2\delta+\epsilon\). There
\[
 h_\eta\le A(2\delta+\epsilon)\le3A\delta,\qquad
 h_\eta p_\eta'\le3A^2<48.
\]
If \(p_\eta'<0\), then
\[
 |h_\eta p_\eta'|\le\frac{\gamma}{I}
 \le\frac{\gamma}{1-8/\gamma}
 \le\gamma+\frac{64}{7}<\eta-54.
\]
The remaining terms in \eqref{eq:sufficient} satisfy
\[
 h_\eta p_\eta/r+\tfrac32p_\eta^2<4+24=28 .
\]
Thus the right side of \eqref{eq:sufficient} is less than \(76\le\eta\)
on the increasing part and less than \(\eta-26\) on the decreasing
part; the case \(p_\eta'=0\) is immediate. This proves
\(f_\eta\in\mathcal A_R(\eta)\) for every parameter pair in the stated
range, without angular sampling.

\subsection{The first correction, uniformly in \texorpdfstring{\(R\)}{R}}

Let \(J_q=\int_1^R r^2q_\eta^2\dd r\).
The inner transition changes mass and energy by \(O(\delta)\),
uniformly in \(R\). For the outer layer, set
\[
 x=\gamma(R-r),\qquad \lambda=\gamma\delta=O(\gamma^{-2}),
 \qquad x_0=\gamma(R-\zeta)\in[\lambda,\lambda+C]\subset[0,4].
\]
At the junction,
\[
 x_0-\lambda=\frac{C}{(R-x_0/\gamma)^2}.
\]
The derivative of \(p_0=C/r^2\) is bounded in absolute value by 6
throughout the wall. Therefore, on this bounded \(x\)-interval,
\begin{equation}\label{eq:uniform-local}
 p_0(R-x/\gamma)=a+O(\gamma^{-1}),\quad
 (R-x/\gamma)^2=R^2+O(\gamma^{-1}),\quad
 x_0=a+O(\gamma^{-1})
\end{equation}
uniformly in \(R\). For example, each of the first two errors is at
most \(24/\gamma\), and \(|x_0-a|\le\lambda+24/\gamma\).
On the outer piece \(q_\eta=(x-\lambda)_+\). Changing the integration
endpoint from \(x_0\) to \(a\) and using
\eqref{eq:uniform-local} gives
\begin{align}
 1-I&=\frac1\gamma\int_0^a(a-x)\dd x+O(\gamma^{-2})
       =\frac{a^2}{2\gamma}+O(\gamma^{-2}),\label{eq:mass-exp}\\
 C-J_q&=\frac{R^2}{\gamma}\int_0^a(a^2-x^2)\dd x+O(\gamma^{-2})
       =\frac{2R^2a^3}{3\gamma}+O(\gamma^{-2}).\label{eq:energy-exp}
\end{align}
For explicit remainder bounds, there are four contributions. The inner
losses are at most \(6\delta\) in mass and \(18\delta\) in energy.
In the outer integral, variation of \(p_0\) contributes at most
\(96/\gamma\) in mass; variation of \(r^2\) and \(p_0^2\) contributes
\((1536+3744)/\gamma\) in energy. The ramp shift obeys
\[
 |(x-\lambda)_+-x|\le\lambda,\qquad
 |(x-\lambda)_+^2-x^2|\le2x\lambda;
\]
its contributions are at most \(4\lambda\) and \(288\lambda\).
Finally, moving the endpoint to \(a\) contributes at most
\(4(\lambda+24/\gamma)\) and \(144(\lambda+24/\gamma)\).
The outer contributions carry the factor \(1/\gamma\). Since
\(\delta\le\gamma^{-3}\), \(\lambda\le\gamma^{-2}\), and \(\gamma\ge64\),
these estimates give
\begin{equation}\label{eq:explicit-remainders}
 \left|I-1+\frac{a^2}{2\gamma}\right|\le\frac{200}{\gamma^2},
 \qquad
 \left|J_q-C+\frac{2K_R}{\gamma}\right|\le\frac{9000}{\gamma^2}.
\end{equation}
Here \(|a-x|\le4\), \(|R^2(a^2-x^2)|\le144\),
\(|p_0^2-(x-\lambda)_+^2|\le16\), and
\(|p_0^2-a^2|\le(3+4/3)24/\gamma=104/\gamma\)
justify the constants above on \(0\le x\le\max(a,x_0)\) where needed.

Using \(I\ge7/8\) and \(C=R^2a\), normalization now gives
\begin{equation}\label{eq:normalized-exp}
 \frac{J_q}{I^2}
 =C+\frac{Ca^2-2R^2a^3/3}{\gamma}+O(\gamma^{-2})
 =C+\frac{K_R}{\gamma}+O(\gamma^{-2}).
\end{equation}
For a numerical remainder constant, put \(t=1-I\le8/\gamma\le1/8\).
Then
\[
 0\le(1-t)^{-2}-1-2t
 =\frac{t^2(3-2t)}{(1-t)^2}\le4t^2,
 \qquad
 \left|I^{-2}-1-\frac{a^2}{\gamma}\right|
 \le\frac{700}{\gamma^2}.
\]
Expanding \(J_q/I^2\) with \eqref{eq:explicit-remainders} bounds the
remainder in \eqref{eq:normalized-exp} by \(14000/\gamma^2\): indeed
\[
 3\cdot700+\frac{32}{9}\frac{16}{9}
 +\frac{32}{9}\frac{700}{64}+9000\frac{64}{49}<14000.
\]
To control smoothing, the translation bound for total variation implies
\[
 \|q_\eta*\kappa_\epsilon-q_\eta\|_{L^1}
 \le\TV(q_\eta)\epsilon\int |t|\kappa(t)\dd t
 \le\frac58 C\epsilon .
\]
Both functions are bounded by \(C\) and supported inside \([1,R]\).
Consequently
\begin{equation}\label{eq:smoothing-error}
 \left|\int_1^R r^2(q_\eta*\kappa_\epsilon)^2\dd r-J_q\right|
 \le \frac54R^2C^2\epsilon=:D_\eta .
\end{equation}
Convolution preserves the mass. Since
\(D_\eta/I^2\le(5/4)81(64/49)\gamma^{-4}\), we obtain the uniform bound
\[
 \left|J(f_\eta)-C-\frac{K_R}{\gamma}\right|
 \le\frac{15000}{\gamma^2}.
\]
Finally, \(\eta/\gamma\le2\) gives
\[
 \eta(J(f_\eta)-C)
 \le K_R+\frac{64K_R}{\gamma}+\frac{15000\eta}{\gamma^2}
 \le K_R+\frac{128K_R+60000}{\eta}
 \le K_R+\frac{61000}{\eta}.
\]
Together with the explicit lower error estimate, this proves
Theorem~\ref{thm:main} with the stated constants.

\section{Rational enclosures and reproduction}\label{sec:numerics}

For rational \(R,\eta\), all construction parameters are rational except
the algebraic junctions. We enclose those junctions by 256 rational
bisection steps in the monotone intervals \eqref{eq:junctions}.
Set \(a_0=1+\delta\), \(b_0=R-\delta\),
\(S=\alpha-a_0\), and \(T=b_0-\zeta\). Exact integration gives
\begin{align}
 I={}&\frac{C}{2\delta}S^2
       +C(1/\alpha-1/\zeta)+\frac{\gamma}{2}T^2,\label{eq:exact-I}\\
 J_q={}&(C/\delta)^2
       \left(\frac{a_0^2S^3}{3}+\frac{a_0S^4}{2}+\frac{S^5}{5}\right)
       +C^2(1/\alpha-1/\zeta)\notag\\
     &+\gamma^2
       \left(\frac{b_0^2T^3}{3}-\frac{b_0T^4}{2}+\frac{T^5}{5}\right).
       \label{eq:exact-J}
\end{align}
Each integral is monotone under enlargement of its interval because its
integrand is nonnegative. Applying this fact separately to the three
pieces produces rational bounds \(I_-\le I\le I_+\) and
\(Q_-\le J_q\le Q_+\). Equation \eqref{eq:smoothing-error} then gives
\begin{equation}\label{eq:rational-upper}
 \frac{Q_--D_\eta}{I_+^2}\le J(f_\eta)
 \le\frac{Q_++D_\eta}{I_-^2},
\end{equation}
with a positive lower numerator in all tabulated cases.
An integer-square-root construction gives a rational upper bound for
\(\sqrt{a/\eta}\), and hence a rational lower bound in
\eqref{eq:finite-lower}.

Table~\ref{tab:enclosures} lists enclosures of
\(\Xi_R(\eta)=\eta[J_R^*(\eta)-C]\).
Its upper endpoint comes from the explicit admissible family, not from
locating an optimizer. All arithmetic preceding decimal presentation
uses exact fractions; decimal endpoints are rounded outward.
The exact rational endpoints and the generating script
\texttt{reproduce.py} accompany this paper.
An additional 120-digit, piecewise quadrature calculation reproduces the
unsmoothed mass and energy inside their rational enclosures.
The finite calculations check the formulas and construction; the
uniform theorem follows from the preceding analytic bounds.

\begin{table}[htbp]
\centering
\small
\caption{Outward-rounded rational enclosures of $\Xi_R(\eta)=\eta[J_R^*(\eta)-C]$. The lower endpoint is $L_R$ from Theorem~\ref{thm:main}; the upper endpoint comes from the explicit admissible family in Section~\ref{sec:recovery}. The last column gives the proved limit.}
\label{tab:enclosures}
\begin{tabular}{rrr r c}
\toprule
$R$ & $\eta$ & Lower endpoint & Upper endpoint & $K_R$ \\
\midrule
$3/2$ & 128 & 1.395307386 & 3.658586574 & $16/9$ \\
$3/2$ & 512 & 1.590796077 & 2.039906892 &  \\
$3/2$ & 2048 & 1.685591794 & 1.836775343 &  \\
$3/2$ & 8192 & 1.732040918 & 1.792168268 &  \\
$3/2$ & 32768 & 1.755002098 & 1.781353645 &  \\
$3/2$ & 131072 & 1.766413588 & 1.778670395 &  \\
\midrule
$2$ & 128 & 0.145264228 & 0.336030470 & $1/6$ \\
$2$ & 512 & 0.156097449 & 0.190686767 &  \\
$2$ & 2048 & 0.161418917 & 0.172082473 &  \\
$2$ & 8192 & 0.164052486 & 0.167988140 &  \\
$2$ & 32768 & 0.165362060 & 0.166995065 &  \\
$2$ & 131072 & 0.166014991 & 0.166748644 &  \\
\midrule
$3$ & 128 & 0.012875360 & 0.028027553 & $1/72$ \\
$3$ & 512 & 0.013384811 & 0.015888803 &  \\
$3$ & 2048 & 0.013637562 & 0.014338348 &  \\
$3$ & 8192 & 0.013763409 & 0.013998455 &  \\
$3$ & 32768 & 0.013826195 & 0.013916111 &  \\
$3$ & 131072 & 0.013857553 & 0.013895684 &  \\
\bottomrule
\end{tabular}
\end{table}


\section{Scope and further questions}

The result concerns the fixed one-component shift in
\eqref{eq:metric}, monotone radial \(C^2\) profiles with exact endpoint
matching, the stated Eulerian normal, and the electric part of the full
Riemann tensor. The wall width does not tend to zero. The theorem neither
bounds the remaining curvature components or other observers' tidal
measurements nor addresses matter realizability, quantum inequalities,
stability, or the satisfaction of energy conditions. Existence and
uniqueness of a finite-\(\eta\) optimizer, sharper remainder estimates,
and additional derivative constraints remain separate questions.

The proof is analytic. The accompanying Lean file records a weaker
statement with existential error constants, rather than the explicit
constants of Theorem~\ref{thm:main}. It formalizes neither the proof nor
the metric-to-matrix derivation. Symbolic algebra and exact rational calculations support
reproducibility. This draft has not undergone human peer review.
The supplementary \texttt{VERIFICATION-RECORD.md} records the computational
and derivational checks.

\paragraph{AI assistance.}
Claude and Codex were used for mathematical exploration, derivation,
implementation, literature organization, and drafting. They are tools,
not authors. The named author is responsible for the content and for
decisions about publication.

\small
\begin{thebibliography}{99}
\bibitem{alcubierre}
Miguel Alcubierre.
\emph{The warp drive: hyper-fast travel within general relativity}.
Classical and Quantum Gravity 11, L73-L77 (1994).
\href{https://doi.org/10.1088/0264-9381/11/5/001}{doi:10.1088/0264-9381/11/5/001}.
\bibitem{pfenningford}
Michael J. Pfenning and L. H. Ford.
\emph{The unphysical nature of warp drive}.
Classical and Quantum Gravity 14, 1743-1751 (1997).
\href{https://doi.org/10.1088/0264-9381/14/7/011}{doi:10.1088/0264-9381/14/7/011}.
\bibitem{lobovisser}
Francisco S. N. Lobo and Matt Visser.
\emph{Fundamental limitations on `warp drive' spacetimes}.
Classical and Quantum Gravity 21, 5871-5892 (2004).
\href{https://doi.org/10.1088/0264-9381/21/24/011}{doi:10.1088/0264-9381/21/24/011}.
\bibitem{bobrickmartire}
Alexey Bobrick and Gianni Martire.
\emph{Introducing physical warp drives}.
Classical and Quantum Gravity 38, 105009 (2021).
\href{https://doi.org/10.1088/1361-6382/abdf6e}{doi:10.1088/1361-6382/abdf6e}.
Also \href{https://arxiv.org/abs/2102.06824v2}{arXiv:2102.06824v2}.
\bibitem{natario}
José Natário.
\emph{Warp drive with zero expansion}.
Classical and Quantum Gravity 19, 1157-1166 (2002).
\href{https://doi.org/10.1088/0264-9381/19/6/308}{doi:10.1088/0264-9381/19/6/308}.
\bibitem{bolivar}
Nelson Bolívar, Ivaylo Vasilev and Gabriel Abellán.
\emph{Zero expansion is not free: the sharp minimum of Eulerian negative energy in Natário warp drives}.
Preprint (2026).
\href{https://arxiv.org/abs/2609.36211v1}{arXiv:2609.36211v1}.
\bibitem{lentz}
Erik W. Lentz.
\emph{Breaking the warp barrier: hyper-fast solitons in Einstein--Maxwell-plasma theory}.
Classical and Quantum Gravity 38, 075015 (2021).
\href{https://doi.org/10.1088/1361-6382/abe692}{doi:10.1088/1361-6382/abe692}.
Also \href{https://arxiv.org/abs/2006.07125v2}{arXiv:2006.07125v2}.
\bibitem{fellheisenberg}
Shaun D. B. Fell and Lavinia Heisenberg.
\emph{Positive energy warp drive from hidden geometric structures}.
Classical and Quantum Gravity 38, 155020 (2021).
\href{https://doi.org/10.1088/1361-6382/ac0e47}{doi:10.1088/1361-6382/ac0e47}.
\bibitem{santiago}
Jessica Santiago, Sebastian Schuster and Matt Visser.
\emph{Generic warp drives violate the null energy condition}.
Physical Review D 105, 064038 (2022).
\href{https://doi.org/10.1103/PhysRevD.105.064038}{doi:10.1103/PhysRevD.105.064038}.
Also \href{https://arxiv.org/abs/2105.03079v2}{arXiv:2105.03079v2}.
\bibitem{warpfactory}
Christopher Helmerich, Jared Fuchs, Alexey Bobrick, Luke Sellers, Brandon Melcher and Gianni Martire.
\emph{Analyzing warp drive spacetimes with Warp Factory}.
Classical and Quantum Gravity 41, 095009 (2024).
\href{https://doi.org/10.1088/1361-6382/ad2e42}{doi:10.1088/1361-6382/ad2e42}.
\bibitem{warpax}
An T. Le.
\emph{Observer-robust energy condition verification for warp drive spacetimes}.
Preprint (2026).
\href{https://arxiv.org/abs/2602.18023v6}{arXiv:2602.18023v6}.
\bibitem{braides}
Andrea Braides and Lev Truskinovsky.
\emph{Asymptotic expansions by $\Gamma$-convergence}.
Continuum Mechanics and Thermodynamics 20, 21-62 (2008).
\href{https://doi.org/10.1007/s00161-008-0072-2}{doi:10.1007/s00161-008-0072-2}.
\end{thebibliography}

\end{document}
